% Options for packages loaded elsewhere
\PassOptionsToPackage{unicode}{hyperref}
\PassOptionsToPackage{hyphens}{url}
\documentclass[
]{article}
\usepackage{xcolor}
\usepackage[margin=1in]{geometry}
\usepackage{amsmath,amssymb}
\setcounter{secnumdepth}{-\maxdimen} % remove section numbering
\usepackage{iftex}
\ifPDFTeX
  \usepackage[T1]{fontenc}
  \usepackage[utf8]{inputenc}
  \usepackage{textcomp} % provide euro and other symbols
\else % if luatex or xetex
  \usepackage{unicode-math} % this also loads fontspec
  \defaultfontfeatures{Scale=MatchLowercase}
  \defaultfontfeatures[\rmfamily]{Ligatures=TeX,Scale=1}
\fi
\usepackage{lmodern}
\ifPDFTeX\else
  % xetex/luatex font selection
\fi
% Use upquote if available, for straight quotes in verbatim environments
\IfFileExists{upquote.sty}{\usepackage{upquote}}{}
\IfFileExists{microtype.sty}{% use microtype if available
  \usepackage[]{microtype}
  \UseMicrotypeSet[protrusion]{basicmath} % disable protrusion for tt fonts
}{}
\makeatletter
\@ifundefined{KOMAClassName}{% if non-KOMA class
  \IfFileExists{parskip.sty}{%
    \usepackage{parskip}
  }{% else
    \setlength{\parindent}{0pt}
    \setlength{\parskip}{6pt plus 2pt minus 1pt}}
}{% if KOMA class
  \KOMAoptions{parskip=half}}
\makeatother
\usepackage{color}
\usepackage{fancyvrb}
\newcommand{\VerbBar}{|}
\newcommand{\VERB}{\Verb[commandchars=\\\{\}]}
\DefineVerbatimEnvironment{Highlighting}{Verbatim}{commandchars=\\\{\}}
% Add ',fontsize=\small' for more characters per line
\usepackage{framed}
\definecolor{shadecolor}{RGB}{248,248,248}
\newenvironment{Shaded}{\begin{snugshade}}{\end{snugshade}}
\newcommand{\AlertTok}[1]{\textcolor[rgb]{0.94,0.16,0.16}{#1}}
\newcommand{\AnnotationTok}[1]{\textcolor[rgb]{0.56,0.35,0.01}{\textbf{\textit{#1}}}}
\newcommand{\AttributeTok}[1]{\textcolor[rgb]{0.13,0.29,0.53}{#1}}
\newcommand{\BaseNTok}[1]{\textcolor[rgb]{0.00,0.00,0.81}{#1}}
\newcommand{\BuiltInTok}[1]{#1}
\newcommand{\CharTok}[1]{\textcolor[rgb]{0.31,0.60,0.02}{#1}}
\newcommand{\CommentTok}[1]{\textcolor[rgb]{0.56,0.35,0.01}{\textit{#1}}}
\newcommand{\CommentVarTok}[1]{\textcolor[rgb]{0.56,0.35,0.01}{\textbf{\textit{#1}}}}
\newcommand{\ConstantTok}[1]{\textcolor[rgb]{0.56,0.35,0.01}{#1}}
\newcommand{\ControlFlowTok}[1]{\textcolor[rgb]{0.13,0.29,0.53}{\textbf{#1}}}
\newcommand{\DataTypeTok}[1]{\textcolor[rgb]{0.13,0.29,0.53}{#1}}
\newcommand{\DecValTok}[1]{\textcolor[rgb]{0.00,0.00,0.81}{#1}}
\newcommand{\DocumentationTok}[1]{\textcolor[rgb]{0.56,0.35,0.01}{\textbf{\textit{#1}}}}
\newcommand{\ErrorTok}[1]{\textcolor[rgb]{0.64,0.00,0.00}{\textbf{#1}}}
\newcommand{\ExtensionTok}[1]{#1}
\newcommand{\FloatTok}[1]{\textcolor[rgb]{0.00,0.00,0.81}{#1}}
\newcommand{\FunctionTok}[1]{\textcolor[rgb]{0.13,0.29,0.53}{\textbf{#1}}}
\newcommand{\ImportTok}[1]{#1}
\newcommand{\InformationTok}[1]{\textcolor[rgb]{0.56,0.35,0.01}{\textbf{\textit{#1}}}}
\newcommand{\KeywordTok}[1]{\textcolor[rgb]{0.13,0.29,0.53}{\textbf{#1}}}
\newcommand{\NormalTok}[1]{#1}
\newcommand{\OperatorTok}[1]{\textcolor[rgb]{0.81,0.36,0.00}{\textbf{#1}}}
\newcommand{\OtherTok}[1]{\textcolor[rgb]{0.56,0.35,0.01}{#1}}
\newcommand{\PreprocessorTok}[1]{\textcolor[rgb]{0.56,0.35,0.01}{\textit{#1}}}
\newcommand{\RegionMarkerTok}[1]{#1}
\newcommand{\SpecialCharTok}[1]{\textcolor[rgb]{0.81,0.36,0.00}{\textbf{#1}}}
\newcommand{\SpecialStringTok}[1]{\textcolor[rgb]{0.31,0.60,0.02}{#1}}
\newcommand{\StringTok}[1]{\textcolor[rgb]{0.31,0.60,0.02}{#1}}
\newcommand{\VariableTok}[1]{\textcolor[rgb]{0.00,0.00,0.00}{#1}}
\newcommand{\VerbatimStringTok}[1]{\textcolor[rgb]{0.31,0.60,0.02}{#1}}
\newcommand{\WarningTok}[1]{\textcolor[rgb]{0.56,0.35,0.01}{\textbf{\textit{#1}}}}
\usepackage{graphicx}
\makeatletter
\newsavebox\pandoc@box
\newcommand*\pandocbounded[1]{% scales image to fit in text height/width
  \sbox\pandoc@box{#1}%
  \Gscale@div\@tempa{\textheight}{\dimexpr\ht\pandoc@box+\dp\pandoc@box\relax}%
  \Gscale@div\@tempb{\linewidth}{\wd\pandoc@box}%
  \ifdim\@tempb\p@<\@tempa\p@\let\@tempa\@tempb\fi% select the smaller of both
  \ifdim\@tempa\p@<\p@\scalebox{\@tempa}{\usebox\pandoc@box}%
  \else\usebox{\pandoc@box}%
  \fi%
}
% Set default figure placement to htbp
\def\fps@figure{htbp}
\makeatother
\setlength{\emergencystretch}{3em} % prevent overfull lines
\providecommand{\tightlist}{%
  \setlength{\itemsep}{0pt}\setlength{\parskip}{0pt}}
\usepackage{bookmark}
\IfFileExists{xurl.sty}{\usepackage{xurl}}{} % add URL line breaks if available
\urlstyle{same}
\hypersetup{
  pdftitle={Continuous distributions and Q-Q plots},
  hidelinks,
  pdfcreator={LaTeX via pandoc}}

\title{Continuous distributions and Q-Q plots}
\author{}
\date{\vspace{-2.5em}}

\begin{document}
\maketitle

\textbf{Inference labs} use the five-step template: Hypothesis →
Visualise → Assumptions → Conduct → Conclude.

\subsection{Learning objectives}\label{learning-objectives}

\begin{itemize}
\tightlist
\item
  Draw densities of the normal, \emph{t}, chi-square, \emph{F}, and
  exponential distributions and describe how shape changes with
  parameters.
\item
  Use Q-Q plots to assess whether data are consistent with a
  hypothesised distribution.
\item
  Recognise the degrees-of-freedom relationships among the
  normal-derived distributions.
\end{itemize}

\subsection{Prerequisites}\label{prerequisites}

Lab 2.4.

\subsection{Background}\label{background}

The continuous distributions in this lab cover the majority of routine
parametric inference. The normal is the distribution of means under the
central limit theorem. The \emph{t} is what happens to the normal when
the variance is itself estimated from the data --- with more uncertainty
(smaller \emph{n}), heavier tails. The chi-square is the distribution of
sums of squared standard normals. The \emph{F} is a ratio of
chi-squares. The exponential appears in survival and queueing.

The Q-Q plot is the workhorse for assessing distributional fit. Plotting
sample quantiles against theoretical quantiles gives a straight line
when the model is right. Deviations at the tails indicate skew or heavy
tails; an S-shape indicates non-normal symmetric behaviour. Q-Q plots
are how assumptions of the \emph{t}-test and the analysis of variance
get checked in practice.

Degrees of freedom are easy to mis-remember. The short version: \emph{t}
has df = \emph{n} − 1 for a one-sample test; chi-square from a sum of
\emph{k} squared standard normals has df = \emph{k}; \emph{F} has two df
--- numerator and denominator --- that come from the chi-squares in its
ratio.

\subsection{Setup}\label{setup}

\begin{Shaded}
\begin{Highlighting}[]
\FunctionTok{library}\NormalTok{(tidyverse)}
\FunctionTok{set.seed}\NormalTok{(}\DecValTok{42}\NormalTok{)}
\FunctionTok{theme\_set}\NormalTok{(}\FunctionTok{theme\_minimal}\NormalTok{(}\AttributeTok{base\_size =} \DecValTok{12}\NormalTok{))}
\end{Highlighting}
\end{Shaded}

\subsection{1. Hypothesis}\label{hypothesis}

We are not testing a hypothesis. We are characterising five continuous
distributions and using Q-Q plots to assess whether a sample is
consistent with a model.

\subsection{2. Visualise}\label{visualise}

Plot densities side by side.

\begin{Shaded}
\begin{Highlighting}[]

\NormalTok{densities }\OtherTok{\textless{}{-}} \FunctionTok{bind\_rows}\NormalTok{(}
\NormalTok{  grid }\SpecialCharTok{|\textgreater{}} \FunctionTok{mutate}\NormalTok{(}\AttributeTok{density =} \FunctionTok{dnorm}\NormalTok{(x),         }\AttributeTok{dist =} \StringTok{"Normal(0,1)"}\NormalTok{),}
\NormalTok{  grid }\SpecialCharTok{|\textgreater{}} \FunctionTok{mutate}\NormalTok{(}\AttributeTok{density =} \FunctionTok{dt}\NormalTok{(x, }\AttributeTok{df =} \DecValTok{3}\NormalTok{),    }\AttributeTok{dist =} \StringTok{"t(3)"}\NormalTok{),}
\NormalTok{  grid }\SpecialCharTok{|\textgreater{}} \FunctionTok{mutate}\NormalTok{(}\AttributeTok{density =} \FunctionTok{dt}\NormalTok{(x, }\AttributeTok{df =} \DecValTok{30}\NormalTok{),   }\AttributeTok{dist =} \StringTok{"t(30)"}\NormalTok{)}
\NormalTok{)}
\end{Highlighting}
\end{Shaded}

\begin{Shaded}
\begin{Highlighting}[]
\FunctionTok{ggplot}\NormalTok{(densities, }\FunctionTok{aes}\NormalTok{(x, density, }\AttributeTok{colour =}\NormalTok{ dist)) }\SpecialCharTok{+}
  \FunctionTok{geom\_line}\NormalTok{(}\AttributeTok{linewidth =} \DecValTok{1}\NormalTok{) }\SpecialCharTok{+}
  \FunctionTok{labs}\NormalTok{(}\AttributeTok{x =} \ConstantTok{NULL}\NormalTok{, }\AttributeTok{y =} \StringTok{"Density"}\NormalTok{, }\AttributeTok{colour =} \ConstantTok{NULL}\NormalTok{)}
\end{Highlighting}
\end{Shaded}

The \emph{t} with 3 df has visibly heavier tails; at 30 df it is nearly
normal.

\begin{Shaded}
\begin{Highlighting}[]
\NormalTok{chi }\OtherTok{\textless{}{-}} \FunctionTok{bind\_rows}\NormalTok{(}
\NormalTok{  chi\_grid }\SpecialCharTok{|\textgreater{}} \FunctionTok{mutate}\NormalTok{(}\AttributeTok{density =} \FunctionTok{dchisq}\NormalTok{(x, }\AttributeTok{df =} \DecValTok{2}\NormalTok{), }\AttributeTok{dist =} \StringTok{"chi{-}sq(2)"}\NormalTok{),}
\NormalTok{  chi\_grid }\SpecialCharTok{|\textgreater{}} \FunctionTok{mutate}\NormalTok{(}\AttributeTok{density =} \FunctionTok{dchisq}\NormalTok{(x, }\AttributeTok{df =} \DecValTok{5}\NormalTok{), }\AttributeTok{dist =} \StringTok{"chi{-}sq(5)"}\NormalTok{),}
\NormalTok{  chi\_grid }\SpecialCharTok{|\textgreater{}} \FunctionTok{mutate}\NormalTok{(}\AttributeTok{density =} \FunctionTok{dchisq}\NormalTok{(x, }\AttributeTok{df =} \DecValTok{10}\NormalTok{), }\AttributeTok{dist =} \StringTok{"chi{-}sq(10)"}\NormalTok{)}
\NormalTok{)}
\end{Highlighting}
\end{Shaded}

\begin{Shaded}
\begin{Highlighting}[]
\FunctionTok{ggplot}\NormalTok{(chi, }\FunctionTok{aes}\NormalTok{(x, density, }\AttributeTok{colour =}\NormalTok{ dist)) }\SpecialCharTok{+}
  \FunctionTok{geom\_line}\NormalTok{(}\AttributeTok{linewidth =} \DecValTok{1}\NormalTok{) }\SpecialCharTok{+}
  \FunctionTok{labs}\NormalTok{(}\AttributeTok{x =} \ConstantTok{NULL}\NormalTok{, }\AttributeTok{y =} \StringTok{"Density"}\NormalTok{, }\AttributeTok{colour =} \ConstantTok{NULL}\NormalTok{)}
\end{Highlighting}
\end{Shaded}

\subsection{3. Assumptions}\label{assumptions}

We assume samples are independent and identically distributed. For the
Q-Q check, the sample must be reasonably sized (say, at least 20
observations) for the plot to be informative.

\begin{Shaded}
\begin{Highlighting}[]
\NormalTok{f\_grid }\OtherTok{\textless{}{-}} \FunctionTok{expand\_grid}\NormalTok{(}
  \AttributeTok{x  =} \FunctionTok{seq}\NormalTok{(}\FloatTok{0.01}\NormalTok{, }\DecValTok{5}\NormalTok{, }\AttributeTok{length.out =} \DecValTok{300}\NormalTok{),}
  \AttributeTok{df1 =} \FunctionTok{c}\NormalTok{(}\DecValTok{2}\NormalTok{, }\DecValTok{5}\NormalTok{, }\DecValTok{10}\NormalTok{),}
  \AttributeTok{df2 =} \FunctionTok{c}\NormalTok{(}\DecValTok{10}\NormalTok{, }\DecValTok{30}\NormalTok{)}
\NormalTok{) }\SpecialCharTok{|\textgreater{}}
  \FunctionTok{mutate}\NormalTok{(}\AttributeTok{density =} \FunctionTok{df}\NormalTok{(x, df1, df2))}
\end{Highlighting}
\end{Shaded}

\begin{Shaded}
\begin{Highlighting}[]
\NormalTok{f\_grid }\SpecialCharTok{|\textgreater{}}
  \FunctionTok{ggplot}\NormalTok{(}\FunctionTok{aes}\NormalTok{(x, density, }\AttributeTok{colour =} \FunctionTok{factor}\NormalTok{(df1))) }\SpecialCharTok{+}
  \FunctionTok{geom\_line}\NormalTok{(}\AttributeTok{linewidth =} \DecValTok{1}\NormalTok{) }\SpecialCharTok{+}
  \FunctionTok{facet\_wrap}\NormalTok{(}\SpecialCharTok{\textasciitilde{}}\NormalTok{ df2, }\AttributeTok{labeller =}\NormalTok{ label\_both) }\SpecialCharTok{+}
  \FunctionTok{labs}\NormalTok{(}\AttributeTok{x =} \ConstantTok{NULL}\NormalTok{, }\AttributeTok{y =} \StringTok{"Density"}\NormalTok{, }\AttributeTok{colour =} \StringTok{"df1"}\NormalTok{)}
\end{Highlighting}
\end{Shaded}

\subsection{4. Conduct}\label{conduct}

Q-Q plots. Simulate data, compare to a hypothesised distribution.

\begin{Shaded}
\begin{Highlighting}[]
\NormalTok{samples }\OtherTok{\textless{}{-}} \FunctionTok{tibble}\NormalTok{(}
  \AttributeTok{normal\_sample =} \FunctionTok{rnorm}\NormalTok{(n),}
  \AttributeTok{heavy\_tail    =} \FunctionTok{rt}\NormalTok{(n, }\AttributeTok{df =} \DecValTok{3}\NormalTok{),}
  \AttributeTok{exponential   =} \FunctionTok{rexp}\NormalTok{(n, }\AttributeTok{rate =} \DecValTok{1}\NormalTok{)}
\NormalTok{)}
\end{Highlighting}
\end{Shaded}

\begin{Shaded}
\begin{Highlighting}[]
\NormalTok{samples }\SpecialCharTok{|\textgreater{}}
  \FunctionTok{pivot\_longer}\NormalTok{(}\FunctionTok{everything}\NormalTok{(), }\AttributeTok{names\_to =} \StringTok{"sample"}\NormalTok{, }\AttributeTok{values\_to =} \StringTok{"value"}\NormalTok{) }\SpecialCharTok{|\textgreater{}}
  \FunctionTok{ggplot}\NormalTok{(}\FunctionTok{aes}\NormalTok{(}\AttributeTok{sample =}\NormalTok{ value)) }\SpecialCharTok{+}
  \FunctionTok{stat\_qq}\NormalTok{(}\AttributeTok{alpha =} \FloatTok{0.6}\NormalTok{) }\SpecialCharTok{+}
  \FunctionTok{stat\_qq\_line}\NormalTok{(}\AttributeTok{colour =} \StringTok{"firebrick"}\NormalTok{) }\SpecialCharTok{+}
  \FunctionTok{facet\_wrap}\NormalTok{(}\SpecialCharTok{\textasciitilde{}}\NormalTok{ sample, }\AttributeTok{scales =} \StringTok{"free"}\NormalTok{) }\SpecialCharTok{+}
  \FunctionTok{labs}\NormalTok{(}\AttributeTok{x =} \StringTok{"Theoretical normal quantiles"}\NormalTok{,}
       \AttributeTok{y =} \StringTok{"Sample quantiles"}\NormalTok{)}
\end{Highlighting}
\end{Shaded}

The normal sample lies on the 45° reference line; the heavy-tailed
sample splays at both ends; the exponential curves away on one side.

Shapiro-Wilk as a quick numerical check on the three samples.

\begin{Shaded}
\begin{Highlighting}[]
  \FunctionTok{pivot\_longer}\NormalTok{(}\FunctionTok{everything}\NormalTok{(), }\AttributeTok{names\_to =} \StringTok{"sample"}\NormalTok{, }\AttributeTok{values\_to =} \StringTok{"value"}\NormalTok{) }\SpecialCharTok{|\textgreater{}}
  \FunctionTok{group\_by}\NormalTok{(sample) }\SpecialCharTok{|\textgreater{}}
  \FunctionTok{summarise}\NormalTok{(}\AttributeTok{p =} \FunctionTok{shapiro.test}\NormalTok{(value)}\SpecialCharTok{$}\NormalTok{p.value, }\AttributeTok{.groups =} \StringTok{"drop"}\NormalTok{)}
\NormalTok{sw\_results}
\end{Highlighting}
\end{Shaded}

\subsection{5. Concluding statement}\label{concluding-statement}

\begin{quote}
The normal sample (\emph{n} = \texttt{n}) was consistent with normality
both visually (Q-Q line) and by Shapiro-Wilk (\emph{p} =
\texttt{signif(sw\_results\$p{[}sw\_results\$sample=="normal\_sample"{]},\ 2)}).
The \emph{t}(3) sample showed clear tail departures, Shapiro-Wilk
\emph{p} =
\texttt{signif(sw\_results\$p{[}sw\_results\$sample=="heavy\_tail"{]},\ 2)}.
The exponential sample was strongly right-skewed, \emph{p} \textless{}
0.001.
\end{quote}

A Q-Q plot is the most information-dense way to inspect a distributional
assumption. Numerical tests of normality (Shapiro-Wilk,
Kolmogorov-Smirnov) should corroborate, not replace, the plot.

\subsection{Common pitfalls}\label{common-pitfalls}

\begin{itemize}
\tightlist
\item
  Reading the centre of a Q-Q plot and ignoring the tails. The tails are
  where the interesting departures live.
\item
  Declaring a distribution wrong on Shapiro-Wilk in a sample of 5000
  when the Q-Q plot is straight.
\item
  Forgetting that \emph{t} approaches normal as df → ∞; do not test with
  df = 500 in a fit-diagnostic mindset.
\end{itemize}

\subsection{Further reading}\label{further-reading}

\begin{itemize}
\tightlist
\item
  Rice JA. \emph{Mathematical Statistics and Data Analysis}.
\item
  Kutner MH et al.~\emph{Applied Linear Statistical Models}.
\end{itemize}

\subsection{Session info}\label{session-info}

\begin{Shaded}
\begin{Highlighting}[]

\end{Highlighting}
\end{Shaded}

\subsection{See also --- chapter index}\label{see-also-chapter-index}

\begin{itemize}
\tightlist
\item
  \href{../index.Rmd}{Methods in this chapter}
\item
  \href{../../../ch00-find-your-method.Rmd}{Find your method (Chapter
  0)}
\end{itemize}

\end{document}
